\documentclass[12pt, a4paper]{article}
\usepackage{../../../myconfig} % Gọi file cấu hình từ ngoài gốc vào

% ==========================================
% BẮT ĐẦU NỘI DUNG TÀI LIỆU
% ==========================================
\begin{document}

% Truyền 3 tham số: {Nhãn cam} {Chữ to chính giữa} {Chữ ở góc phải Header}
\titlebox{Chủ đề: Hàm số}{BTDD tốc độ thay đổi một đại lượng}{Hàm số: Tốc độ thay đổi đại lượng}

% --- CÂU 1---
\questionTrueFalse{0}{1}{Sau \( x \) tháng kể từ khi hãng xe buýt STER được khởi chạy tại thành phố MT, số lượng người sử dụng xe buýt này được biểu thị bởi hàm số  
\(N(x) = 4x^3 + 250x + 6000\) 
(người). Biết rằng xe buýt được khởi chạy vào ngày 01/01/2026.}{
    \textbf{a)} & Theo thời gian, số lượng người dùng xe buýt STER tại thành phố MT luôn tăng.
    & \textcolor{amberorange}{$\bigcirc$} & \textcolor{amberorange}{$\bigcirc$} \\ \hline
    
    \textbf{b)} & Trung bình số lượng người dùng xe buýt STER tại thành phố MT trong một tháng đạt giá trị nhỏ nhất bằng 1240 người/tháng (kết quả làm tròn đến hàng đơn vị).
    & \textcolor{amberorange}{$\bigcirc$} & \textcolor{amberorange}{$\bigcirc$} \\ \hline
    
    \textbf{c)} & Trung bình số lượng người dùng xe buýt STER tại thành phố MT trong một tháng đạt giá trị nhỏ nhất vào ngày 03/10/2026.
    & \textcolor{amberorange}{$\bigcirc$} & \textcolor{amberorange}{$\bigcirc$} \\ \hline
    
    \textbf{d)} & Tốc độ thay đổi trung bình số lượng người dùng xe buýt STER tại thành phố MT trong một tháng luôn tăng.
    & \textcolor{amberorange}{$\bigcirc$} & \textcolor{amberorange}{$\bigcirc$} \\
}

% --- CÂU 10 ---
\questionTrueFalse{0}{2}{Số lượng cá thể vi khuẩn của một quần thể sau \( t \) giờ kể từ khi tiêm một loại thuốc được cho bởi công thức  
\(P(t) = \dfrac{48t + 20}{t^2 + 1}\) 
(triệu con vi khuẩn).}{
    \textbf{a)} & Sau 2 giờ kể từ khi tiêm thuốc, quần thể này có 23,2 triệu con vi khuẩn.
    & \textcolor{amberorange}{$\bigcirc$} & \textcolor{amberorange}{$\bigcirc$} \\ \hline
    
    \textbf{b)} & Theo thời gian dài (có thể coi như tiến đến vô cùng), số lượng cá thể vi khuẩn của quần thể này sẽ tiến đến 0.
    & \textcolor{amberorange}{$\bigcirc$} & \textcolor{amberorange}{$\bigcirc$} \\ \hline
    
    \textbf{c)} & Sau 45 phút kể từ khi tiêm thuốc, số lượng cá thể vi khuẩn trong quần thể này đạt giá trị lớn nhất.
    & \textcolor{amberorange}{$\bigcirc$} & \textcolor{amberorange}{$\bigcirc$} \\ \hline
    
    \textbf{d)} & Sau 78 phút kể từ khi tiêm thuốc, tốc độ thay đổi số lượng cá thể vi khuẩn trong quần thể này đạt giá trị nhỏ nhất (kết quả làm tròn đến hàng đơn vị).
    & \textcolor{amberorange}{$\bigcirc$} & \textcolor{amberorange}{$\bigcirc$} \\
}

% --- CÂU 20 ---
\questionShortAnswer{0}{3}{Vào đầu tháng 4/2016, Công ty F đã xả thải ra biển miền T, gây ra sự cố ô nhiễm môi trường nghiêm trọng, làm cá chết hàng loạt dọc bờ biển bốn tỉnh HT, QH, MT, và ST. Nguyên nhân được xác định là do công ty đã xả thải chưa qua xử lý, có chứa độc tố phenol, xyanua ra môi trường. Sự cố này đã gây ra những hậu quả nặng nề cho môi trường và người dân địa phương, đặc biệt là ngành thủy sản.

Công ty F muốn chi ra \( C \) triệu đồng để xử lý \( p\% \) lượng hóa chất có trong nước thải của nhà máy này, với biểu thức liên hệ giữa \( C \) và \( p \) là \( C(p) = \dfrac{310p}{(10 + p)(100 - p)} \). Tại giá trị \( p \) mà tốc độ tăng số tiền cần chi cho việc xử lý hóa chất đạt giá trị nhỏ nhất, công ty cần chi ra bao nhiêu triệu đồng để xử lý hóa chất (kết quả làm tròn đến hàng phần trăm)?}

% --- CÂU 21 ---
\questionShortAnswer{0}{4}{Trong một thử nghiệm về việc học tiếng Anh, thời gian cần thiết để một bạn học sinh khá học được phiên âm của \( n \) từ vựng mới được mô hình hóa bởi hàm số \( f(n) = 4n\sqrt{n-4} \) (đơn vị: giây). Khi tốc độ học phiên âm của bạn học sinh đạt mức thấp nhất, thì khi đó bạn học sinh sẽ học được phiên âm của một từ vựng trong trung bình là bao nhiêu giây?}

% --- CÂU 23 ---
\questionShortAnswer{0}{5}{Tại cửa hàng thời trang STER Shop, chủ STER muốn theo dõi hiệu suất làm việc của các bạn nhân viên đóng gói của cửa hàng. Ông chủ thấy rằng sau \( x \) giờ kể từ khi vào ca làm việc, một nhân viên đóng gói sẽ đóng xong được \( f(x) = -x^3 + 6x^2 + 15x \) sản phẩm. Ngoài ra, trong ca làm việc, ông chủ sẽ cho phép các bạn nhân viên nghỉ giải lao 15 phút, và ông thấy rằng sau \( x \) giờ kể từ khi kết thúc giải lao, một nhân viên đóng gói sẽ đóng xong được \( g(x) = -\dfrac{1}{3}x^3 + x^2 + 23x \) sản phẩm. Biết rằng ca làm việc bắt đầu vào lúc 13 giờ và kết thúc lúc 17 giờ 45 phút, và trong ca làm sẽ có 1 lần nghỉ giải lao 15 phút. Hỏi để các bạn nhân viên đóng gói được nhiều sản phẩm nhất khi ca làm kết thúc, thời gian nghỉ giải lao cần bắt đầu sau bao nhiêu giờ kể từ khi bắt đầu ca làm việc (kết quả làm tròn đến hàng phần trăm)?}





\end{document}